\BeleskeID{04-s028}
% element: 04-s028:e01 — Pitanje jednoznačnosti i negativnih brojeva
% element: 04-s028:e02 — Prelazak −Vp=−(Vp)=−(Vq)=−Vq preko apsolutne vrednosti
% element: 04-s028:e03 — Digitalni sistem nema poseban nivo za minus; potrebno kodovanje
% element: 04-s028:e04 — Četvorobitni sistem, 16 kodova, dve polovine i celobrojne vrednosti
% element: 04-s028:d01

Matematički zapis $-V$ koristi poseban znak minus ispred apsolutne vrednosti. Promena osnove apsolutne vrednosti ne menja taj znak. U binarnom registru imamo samo nule i jedinice, pa i znak mora da bude uključen u dogovor o značenju kodne reči.

Četiri bita daju $2^4=16$ kodova ukupno. Rezervisanje dela kodova za negativne vrednosti smanjuje opseg nenegativnih vrednosti u poređenju sa neoznačenim prikazom. Kako se raspoređuju kodovi, da li nula ima dva prikaza i koje su krajnje vrednosti, zavisi od izabranog kodovanja.
